\section{Finite acquisition: centroid cancellation and cell ranges}\label{sec:acquisition}

The preceding projection admits two rigorous finite acquisition methods.
The first uses counts and first moments, and is insensitive to internal
root gaps. The second uses only counts and certified cell geometry, and
can exploit Gaussian decay when cells are separated. Their inputs differ;
neither method is uniformly superior. We derive both error bounds before
giving the finite census as supporting evidence in
Appendix~\ref{app:status}. The source arguments are
\cite{OCCentroids,OCCellRanges}.

\subsection{A complete theta tail with explicit constants}

We first pay for every zero outside the retained rectangle, without
requiring a numerical value of the coarse constant $A$.
The source majorant \eqref{eq:source-basic-majorant} gives
$\Phi(u)<10\ee^u r_u^2\ee^{-r_u}$: consecutive terms in
$\sum_{l\ge1}l^4\ee^{-\pi(l^2-1)}$ have ratio at most
$16\ee^{-3\pi}<1/2$.
Since $u^2\le r_u/12$, substitution $v=r_u$ gives
\begin{equation}\label{eq:explicit-source-majorant}
 \mathcal M:=\int_0^\infty\ee^{u^2}\Phi(u)\dd u
 \le\frac{10}{4\pi^{1/4}}\int_\pi^\infty v^{5/4}\ee^{-11v/12}\dd v
 <5(12/11)^3<7.
\end{equation}
The inequality $u^2\le r_u/12$ follows from $\ee^{2u}\ge2u$ and
$\pi>3$; for the last estimate use $v^{5/4}\le v^2$ and extend the
integral to zero. On $0\le u\le1/10$, $3<r_u<8$, and the first
positive source term is at least $9\ee^{-8}$. Thus
$H_0(0)>1/7290$ and $\log(\mathcal M/H_0(0))<20$.

For $|z|\le2r$, $|\cos(zu)|\le\ee^{2ru}\le\ee^{r^2+u^2}$.
Jensen's formula on the disk of radius $2r$ now gives
\begin{equation}\label{eq:global-count-jensen}
 N_{\rm disk}(r)\log2\le r^2+20,
 \qquad N_{\rm disk}(r)\le2r^2+40.
\end{equation}
Boundary zeros follow by a larger-radius limit. Each real unit slab
starting at $q$ lies, on the canonical strip, in a disk of radius
$|q|+2$; it therefore contains at most $16(|q|+2)^2$ labels.

For $T=1000$ retain
$\mathcal R=(988,2012)+\ii(-3/2,3/2)$.
An omitted label has $|x_j-X|\ge12$ at every center.
Two unit slabs per shell bound its total Gaussian height weight by
\begin{equation}\label{eq:explicit-weight-tail}
 b=32\sum_{n=12}^\infty(2003+n)^2\ee^{-n^2/2}<2^{-44}.
\end{equation}
To verify the strict bound, write $n=12+j$ and dominate the summand by
$2015^2(1+j)^2\,2^{-72}2^{-12j}$. For $q=2^{-12}$,
$\sum_{j\ge0}(1+j)^2q^j=(1+q)/(1-q)^3<2$ and $2015<2^{11}$.
If the complete retained count is $N\le1024$, then
\begin{equation}\label{eq:explicit-projection-tail}
 0\le P_{1000}-P_{1000}^{\mathcal R}
       \le2c_G(N+b)b<2^{-32}.
\end{equation}
This includes negative heights and both infinite horizontal tails.
The count cap is a stated input; the Gaussian decay does not certify
the retained census by itself.

\subsection{Counts and first moments give a second-order error}

On $\mathcal R$, $Z(z)=\zeta(1/2-\ii z/2)$ has the same divisor as
$H_0$ and has no pole or trivial zero. Its completion factors are
holomorphic and nonzero there. On a zero-free cell boundary,
the argument principle and its first-moment version give
\begin{equation}\label{eq:contour-centroid}
 n_C=\frac1{2\pi\ii}\oint_{\partial C}\frac{Z'}Z\dd z,
 \qquad q_C=\frac1{2\pi\ii}\oint_{\partial C}(z-a_C)\frac{Z'}Z\dd z,
 \qquad c_C=a_C+q_C/n_C.
\end{equation}
Residues show that $n_C$ counts labels and $c_C$ is their exact
unweighted complex centroid. A zero-free completion multiplier
contributes zero to both contour integrals. These are classical
contour identities; the useful issue is the uncertainty they remove
from the specified pair projection.

\begin{lemma}[Kernel derivative bounds]\label{lem:kernel-derivatives}
For the smooth real extension of \eqref{eq:projection-kernel} to
$\C^2$, with real two-coordinate derivatives,
\[
 \|\nabla_zK\|,\|\nabla_wK\|\le9/T,\quad
 \|D_z^2K\|_{\rm op},\|D_w^2K\|_{\rm op}\le18/T,
 \quad\|\nabla D\|\le4/T,\quad\|D^2D\|_{\rm op}\le10/T.
\]
\end{lemma}
\begin{proof}
Before center integration the kernel is
\[
 \frac{c_G}{T}\int_T^{2T}w(x-X)w(x'-X)g(x-x')b_+(y)b_-(y')\dd X,
 \quad w(s)=\ee^{-s^2/2},\quad g(d)=\ee^{-4d^2/9}.
\]
These four Gaussian factors and their first two derivatives have
absolute value at most one. In $z$ derivatives, keep $w(x'-X)$
and integrate it over the whole line, with mass $\sqrt{2\pi}<3$.
Gradient components are bounded by $6/T,3/T$, and Hessian entries
in $xx,xy,yy$ by $12/T,6/T,3/T$. The Euclidean gradient and
maximum absolute row sum give $9/T,18/T$. The other variable is identical.
For $D$, use $|I_T|,|I_T'|\le2/T$, $|I_T''|\le8/T$,
and derivative bounds $1,2$ for $\ee^{-16y^2/17}$.
The Hessian entries are bounded by $8/T,2/T,4/T$, proving the claim.
\end{proof}

\begin{theorem}[Centroid acquisition error]\label{thm:centroid}
Partition the retained divisor into convex cells with counts $n_C$,
centroids $c_C$ and total count $N$. Set
$S_2=\sum_C\sum_{z_j\in C}|z_j-c_C|^2$ and
\[
 P_{\rm bar}=\sum_{C,E}n_Cn_EK(c_C,c_E)-\sum_Cn_CD(c_C).
\]
Then
\begin{equation}\label{eq:centroid-error}
 |P_T^{\mathcal R}-P_{\rm bar}|\le\frac{18N+5}{T}S_2.
\end{equation}
In particular cells of diameter at most $r$ give
$S_2\le Nr^2$. If $|\widehat c_C-c_C|\le\sigma$, replacing
centroids by these approximations costs at most
$(18N^2+4N)\sigma/T$.
\end{theorem}
\begin{proof}
For a real $C^2$ test $F$, Taylor expansion in a cell gives
\[
 \left|\sum_{z_j\in C}F(z_j)-n_CF(c_C)\right|
 \le\tfrac12\sup_C\|D^2F\|_{\rm op}
                         \sum_{z_j\in C}|z_j-c_C|^2,
\]
since $\sum_{z_j\in C}(z_j-c_C)=0$ cancels both real components
of the linear term. Apply this successively to both arguments of the
raw all-label sum; each application costs $9NS_2/T$. The self term
costs $5S_2/T$. Convexity puts the centroid inside its cell, so
$S_2\le Nr^2$. Gradient bounds applied to both arguments and then
the diagonal prove the rounding allowance.
\end{proof}
No simple-root or gap hypothesis enters this theorem. In a collapsed
cell the factorial coefficient is $n_C(n_C-1)$, and the cost of
collapsing distinct roots is already paid by \eqref{eq:centroid-error}.
The weights stay inside $K$; replacing these unweighted centroids
by an unexplained weighted average would lose the cancellation.

\subsection{Direct ranges exploit separation and Gaussian decay}

Suppose instead that each of $N$ rectangular cells contains exactly one
label, and that they are sorted horizontally with adjacent interval
distance at least $s>0$. Let $u$ and $v$ bound their horizontal and
vertical widths. For $C<E$ form the exact intervals
\[
 M_{CE}=[(x_C^-+x_E^-)/2,(x_C^++x_E^+)/2]=[m_-,m_+],
 \qquad D_{CE}=[x_E^--x_C^+,x_E^+-x_C^-]=[d_-,d_+].
\]
Here $d_->0$. Define ideal, exact-real endpoints
\begin{align}
 A_-&=\frac{\sqrt\pi}{2T}
   [\operatorname{erf}(2T-m_+)-\operatorname{erf}(T-m_-)],\nonumber\\
 A_+&=\frac{\sqrt\pi}{2T}
   [\operatorname{erf}(2T-m_-)-\operatorname{erf}(T-m_+)],\nonumber\\
 I_-&=\max(0,A_-),\quad I_+=\min(2/T,A_+),\qquad
 G_-=\ee^{-25d_+^2/36},\quad G_+=\ee^{-25d_-^2/36}.
                                                        \label{eq:cell-endpoints}
\end{align}
Monotonicity of $\operatorname{erf}$ gives the enclosure of $I_T$;
clipping uses $0\le I_T\le\sqrt\pi/T<2/T$.
For each mark let $s_{C,\pm}^-,s_{C,\pm}^+$ be the minimum and maximum
of $(y\mp1/2)^2$ on the cell's vertical interval, and put
\[
 B_{C,\pm}^-=\ee^{-8s_{C,\pm}^+/17},\qquad
 B_{C,\pm}^+=\ee^{-8s_{C,\pm}^-/17},\qquad
 H_{CE}^\pm=B_{C,+}^\pm B_{E,-}^\pm+B_{E,+}^\pm B_{C,-}^\pm.
\]
The last expression sums both ordered mark products. Set
\begin{equation}\label{eq:ideal-cell-sum}
 L_*=c_G\sum_{C<E}I_-G_-H_{CE}^-,\qquad
 U_*=c_G\sum_{C<E}I_+G_+H_{CE}^+.
\end{equation}
All factors are nonnegative, so $P_T^{\mathcal R}\in[L_*,U_*]$.

\begin{theorem}[Gaussian range-width bound]\label{thm:cell-ranges}
Under these singleton-cell and separation hypotheses, the ideal radius obeys
\begin{equation}\label{eq:range-width}
 \frac{U_*-L_*}{2}
   \le\frac{N(10u+4v)}{T}\sum_{k=1}^{N-1}\ee^{-(ks)^2/4}
   <\frac{N(10u+4v)\sqrt\pi}{sT}
\end{equation}
when $N\ge2$ and $u+v>0$. For $N\le1$ the retained factorial sum is zero.
\end{theorem}
\begin{proof}
At index separation $k$, every admissible horizontal gap is at least
$ks$. Put $q_k=\ee^{-(ks)^2/4}$. Endpoint differentiation gives
$\operatorname{wid}I\le2u/T$ and $\sup I\le2/T$.
For $g(d)=\ee^{-25d^2/36}$,
\[
 g(d)\le\ee^{-d^2/4},\qquad |g'(d)|\le2\ee^{-d^2/4}.
\]
Indeed write $a=25/36$, use
$|g'|=2\sqrt a(\sqrt a|d|\ee^{-ad^2/2})\ee^{-ad^2/2}$,
$\sup_{t\ge0}t\ee^{-t^2/2}=\ee^{-1/2}$ and $a/2>1/4$.
Since the gap interval has width at most $2u$,
$\sup G\le q_k$ and $\operatorname{wid}G\le4u q_k$.
Each mark is at most one and has derivative at most
$\sqrt{16/(17\ee)}<1$, so its interval width is at most $v$.
Telescoping the product of nonnegative intervals shows that the width
of one ordered contribution is at most
\[
 \frac{2u}{T}q_k+\frac2T\,4u q_k+\frac2Tq_k(2v)
       =\frac{10u+4v}{T}q_k.
\]
There are $2(N-k)$ ordered pairs of this index separation. Summing
and dividing by two proves the first bound. The decreasing Gaussian
has sum strictly below its integral from zero to infinity,
$\sqrt\pi/s$, proving the second.
\end{proof}
The arbitrary-$s$ form makes explicit the geometric input to the
specialized estimate in \cite{OCCellRanges}. It is a sufficient error
bound for separated singleton cells, not a uniform canonical gap theorem.
For general counts, cross-cell coefficients become $n_Cn_E$;
same-cell coefficients $n_C(n_C-1)$ require an enclosure of $K$ on
$C\times C$, not only the diagonal $D$. Such cells do not inherit
\eqref{eq:range-width} without a new width analysis.

\begin{corollary}[Separate arithmetic and geometric allowances]\label{cor:range-arithmetic}
Suppose enclosing balls for the exact endpoints $L_*,U_*$ have radii
$r_L,r_U$, and let $a$ be the lower endpoint of the first ball and $b$
the upper endpoint of the second. If the omitted tail is at most $t$,
then $[a-t,b+t]$ encloses the complete projection and has radius at most
\begin{equation}\label{eq:range-total}
 \frac{U_*-L_*}{2}+r_L+r_U+t.
\end{equation}
\end{corollary}
\begin{proof}
A ball containing an exact endpoint can place its outward edge as much
as twice its radius away from that endpoint. Thus
$b-a\le U_*-L_*+2r_U+2r_L$. The symmetric tail enlargement adds $t$
to the radius. The nonnegative tail would also allow the sharper
asymmetric enlargement $[a,b+t]$; we retain the conservative formula.
\end{proof}
Endpoint arithmetic, the certified census, and the analytic tail are
three distinct premises. A width theorem does not evaluate either
endpoint. The supporting example below retains this distinction.
